% !TEX root = 00_main_aaai_2027.tex
\section{What Scalar Criteria Identify}
\label{sec:problem}
Let $p(r\mid x)$ be the distribution of a specified rating population and let a judge report $q(r\mid x)$.
Write $\mu_q(x)=\sum_r r q(r\mid x)$.
Pearson correlation measures linear association between scalar summaries; Spearman and Kendall measure ordering with their respective tie conventions.
A strictly increasing transformation preserves ranks, but compression into two bins generally introduces ties and can change rank correlation.
None of these statistics identifies the full conditional distribution when different forecasts share $\mu_q$.

Squared error also has a narrower target.
For any scalar prediction $a(x)$,
\begin{align}
\mathbb{E}[(a-R)^2\mid x]
&=(a-\mathbb{E}[R\mid x])^2 \\
&\quad+\operatorname{Var}(R\mid x).
\end{align}
It identifies the conditional mean under finite second moments, not an arbitrary conditional distribution.
Scaling a bounded rating to $[0,1]$ changes the units, not that identification property.

These limitations do not establish that every psychometric method is unnecessary.
Construct validity concerns the interpretation supported by a body of evidence, rather than a single agreement coefficient~\citep{cronbach1955construct}.
Choosing whose judgments to model and whether the rating task measures the intended construct remain necessary even with a proper probability loss.
